\documentclass[11pt,a4paper]{article}
\usepackage[margin=2.5cm]{geometry}
\usepackage{booktabs}
\usepackage{tabularx}
\usepackage{array}
\usepackage{parskip}
\usepackage{helvet}
\renewcommand{\familydefault}{\sfdefault}
\usepackage{microtype}
\usepackage{xcolor}
\definecolor{uniblue}{RGB}{0,83,159}

\begin{document}

\begin{center}
  {\Large\bfseries\color{uniblue} Pre-Study Parameter Configuration}\\[4pt]

\end{center}

\medskip
\noindent\rule{\linewidth}{0.4pt}
\medskip

The selected NPR shader combines multiple rendering techniques in a single pass: an inverted hull geometry outline, Gaussian pre-filtered Sobel edge detection, XToon 2D ramp tonal shading, and rim light. Post-process effects including a Kuwahara painterly filter and a hierarchical depth edge layer can be added on top.

Although the shader includes a Gaussian pre-filter stage that smooths the luminance signal before Sobel edge detection, this stage was disabled across all pre-study configurations. Preliminary testing showed no perceptually meaningful visual difference between the pre-filter enabled and disabled at the parameter ranges used for this avatar, so a dedicated round for this parameter was omitted.

The study consists of five rounds. The rim light round was removed to reduce participant burden. The two separate Kuwahara rounds were merged into a single round with kernel size fixed at 8 and hardness/sharpness varied. Each round varied one parameter group while all others were held constant. Participants ranked images by trustworthiness as an advisor. Images are labelled by round and letter (e.g.\ R2 C = Round~2, image~C).

\bigskip

%% ROUND 1
\noindent{\bfseries Round 1 \textemdash\ Inner Line Threshold (R1 A\textendash E)}

\smallskip
\begin{tabular}{ll}
\toprule
\textbf{Label} & \textbf{\texttt{\_InnerLineThreshold}} \\
\midrule
R1 A & 0.01 \\
R1 B & 0.02 \\
R1 C & 0.03 \\
R1 D & 0.04 \\
R1 E & 0.05 \\
\bottomrule
\end{tabular}

\smallskip
\noindent\textit{Fixed: Gaussian blur off, normal edges off, inner line sample distance 0.2, light sensitivity 0.5, shadow strength 0.5, outline width 0.005, depth offset on, rim light on.}

\bigskip

%% ROUND 2
\noindent{\bfseries Round 2 \textemdash\ Outer Outline Width and Depth Offset (R2 A\textendash F)}

\smallskip
\begin{tabular}{llll}
\toprule
\textbf{Label} & \textbf{Outline Width} & \textbf{Depth Offset} & \textbf{Depth Bias} \\
\midrule
R2 A & 0.005 & Off & 15 \\
R2 B & 0.005 & On  & 15 \\
R2 C & 0.007 & Off & 15 \\
R2 D & 0.007 & On  & 15 \\
R2 E & 0.003 & Off & 15 \\
R2 F & 0.003 & On  & 15 \\
\bottomrule
\end{tabular}

\smallskip
\noindent\textit{Fixed: inner line threshold 0.05 (R1 E). Depth bias 15 throughout to make on/off difference clearly visible (effect is negligible below 5).}

\bigskip

%% ROUND 3
\noindent{\bfseries Round 3 \textemdash\ XToon Toon Shading (R3 A\textendash E)}

\noindent Detail mode fixed to Depth. Curvature and Manual excluded (near-identical output at 1\,m framing).

\smallskip
\begin{tabular}{lll}
\toprule
\textbf{Label} & \textbf{Shadow Strength} & \textbf{Detail Bias} \\
\midrule
R3 A & 1.0 & 0.5 \\
R3 B & 0.0 & 0.5 \\
R3 C & 0.5 & 0.5 \\
R3 D & 0.5 & 1.0 \\
R3 E & 0.5 & 0.0 \\
\bottomrule
\end{tabular}

\smallskip
\noindent\textit{Fixed: outline width 0.005, depth offset on (bias 15), inner line threshold 0.03.}

\bigskip

%% ROUND 4
\noindent{\bfseries Round 4 \textemdash\ Kuwahara Hardness and Sharpness (R4 A\textendash D)}

\noindent Kuwahara applied on top of V4 base (R1 E). Kernel size fixed at 8 (no visible difference observed above this value). R4 A is the V4 baseline without Kuwahara. Hardness and sharpness are varied inversely across three combinations. Sector count excluded (no perceptible difference across its range).

\smallskip
\begin{tabular}{lll}
\toprule
\textbf{Label} & \textbf{Hardness} & \textbf{Sharpness} \\
\midrule
R4 A & --- & --- \\
R4 B & 18  & 1   \\
R4 C & 9   & 9   \\
R4 D & 1   & 18  \\
\bottomrule
\end{tabular}

\smallskip
\noindent\textit{Kernel size: 8 for all Kuwahara variants.}

\bigskip

%% ROUND 5
\noindent{\bfseries Round 5 \textemdash\ Hierarchical Depth Edge Threshold (R5 A\textendash F)}

\noindent Only depth edges tested. Colour edges already handled by V4 Sobel. Normal edges excluded (camera-angle-dependent, unsuitable for fixed-view ranking). Depth scale fixed at 2. R5 A is the V4 baseline without hierarchical depth edges.

\smallskip
\begin{tabular}{ll}
\toprule
\textbf{Label} & \textbf{Depth Threshold} \\
\midrule
R5 A & --- (base, no hierarchical depth edge) \\
R5 B & 0.0  \\
R5 C & 0.02 \\
R5 D & 0.03 \\
R5 E & 0.04 \\
R5 F & 0.05 \\
\bottomrule
\end{tabular}

\bigskip
\noindent\rule{\linewidth}{0.4pt}

\noindent{\bfseries Baseline V4 Configuration (R1 E, used as reference in Rounds 4\textendash 5)}

\smallskip
\begin{tabular}{ll}
\toprule
\textbf{Parameter} & \textbf{Value} \\
\midrule
Inner line threshold        & 0.05 \\
Inner line sample distance  & 0.2  \\
Gaussian blur               & Off  \\
Normal edges                & Off  \\
Fresnel edges               & Off  \\
Light sensitivity           & 0.5  \\
Shadow strength             & 0.5  \\
Detail mode                 & Depth \\
Detail bias                 & 0.5  \\
Outer outline width         & 0.005 \\
Depth offset                & On   \\
Depth bias                  & 15   \\
Rim light                   & On   \\
Rim power                   & 4.0  \\
\bottomrule
\end{tabular}

\end{document}
